In dieser Arbeit diskutieren wir ein mikroskopisches Gittermodell, für das in \cite{EP_DynamicInduced_Quenched_Lehmann_2021} bereits mittels eines vollständig numerischen Ansatzes das Auftreten von sogenannten Exceptional Points (EPs) gezeigt wurde.
Durch Anwendung von Bosonisierung erhalten wir ein hermitesches Luttinger Liquid Modell und identifizieren einen sine-Gordon-artigen Wechselwirkungsterm, der für die EPs verantwortlich ist.
Um die Existenz einer nicht-hermiteschen Topologie analytisch nachzuweisen, berechnen wir die Ein-Teilchen Greenfunktion und wandeln sie in einen effektiven Hamiltonian um.
Da eine exakte Berechnung nicht realistisch ist, approximieren wir das Ergebnis stattdessen mithilfe von Störungstheorie bis zur ersten Ordnung in der Wechselwirkung.
Wir zeigen, dass endliche Temperatur und ein nichttrivialer Luttinger Parameter $K\neq 1$ für die EPs notwendig sind, und dass die Ergebnisse qualitativ mit der Numerik übereinstimmen.
In bestimmten Parameterbereichen ist es au{\ss}erdem möglich, sogar eine quantitative Übereinstimmung zu erreichen, indem Modellparameter an die Daten angepasst werden.
%Wir zeigen, dass unser analytisches Ergebnis auch quantitativ mit der Numerik in bestimmten Parameterbereichen übereinstimmen kann, indem wir Modellparameter an die Daten anpassen.
Abschlie{\ss}end diskutieren wir, wie diese optimalen Werte hypothetisch mithilfe einer Renormierungsgruppe vorhergesagt werden könnten.
Um diese Idee zu veranschaulichen, wenden wir zwei verschiedene, auf Störungstheorie basierende Methoden auf unser Modell an und vergleichen ihre Ergebnisse.
\nextabstract[english]
In this thesis, we discuss a microscopic lattice model, for which the emergence of \glspl{ep} has already been shown in \cite{EP_DynamicInduced_Quenched_Lehmann_2021} via a fully numeric approach.
By applying bosonization, we obtain a Hermitian \acrlong{ll} model and identify the novel \acrlong{sg}-like interaction term responsible for the \glspl{ep}.
To demonstrate the existence of \acrlong{nh} topology by analytical means, we compute the single-particle \acrlong{gf} and convert it into an effective Hamiltonian. 
Since an exact calculation is infeasible, we instead approximate the result using perturbation theory to first order in the interaction.
%The unusual structure of this term leads to subtle analytical problems and requires
We show that finite temperature and a non-trivial Luttinger parameter $K\neq 1$ are required for the formation of \glspl{ep}, and that the results are in qualitative agreement with the numerics.
In certain parameter regimes, it is even possible to quantitatively match the data by tuning model parameters.
%Furthermore our analytic result can also quantitatively match the numerics in certain parameter regimes by using model parameters to fit to the data.
Finally, we discuss how these optimal values may hypothetically be predicted using a \acrlong{rg} approach.
To illustrate this idea, we apply two different perturbative schemes to our model and compare their results.
\glsresetall